\section{自动微分与框架实现}

\subsection{自动微分}

既然反向传播算法是如此系统化，能否让计算机\textbf{完全自动}完成？

\begin{knowledgebox}{自动微分的历史}
早期的深度学习框架 \textbf{Theano}（蒙特利尔大学开发）尝试了完全符号化的自动微分：给定前向计算的符号表达式，自动推导出反向传播的符号表达式。但这种方式过于重量级，难以灵活处理各种情况。
\end{knowledgebox}

现代深度学习框架（PyTorch、TensorFlow 等）采用了一种折中方案：

\begin{itemize}
    \item \textbf{框架负责}：管理计算图、执行拓扑排序、运行前向/反向遍历、传递上游梯度
    \item \textbf{用户负责}：为每种操作实现 \texttt{forward}（前向计算）和 \texttt{backward}（局部梯度计算）方法
\end{itemize}

\subsection{Forward/Backward API}

\begin{figure}[H]
\centering
\includegraphics[width=0.95\textwidth]{slides-images/slide-080.jpg}
\caption{Forward/Backward API 示例：\texttt{MultiplyGate} 类实现 \texttt{forward}（计算 $z = xy$）和 \texttt{backward}（计算 $\frac{\partial L}{\partial x}, \frac{\partial L}{\partial y}$）\protect\footnotemark}
\end{figure}
\footnotetext{Slides 第81页。}

\begin{lstlisting}[caption={乘法门的 Forward/Backward 实现},language=Python]
class MultiplyGate(object):
    def forward(self, x, y):
        z = x * y
        self.x = x  # 缓存输入，供 backward 使用
        self.y = y
        return z

    def backward(self, dz):
        dx = dz * self.y  # 上游梯度 * 对方输入
        dy = dz * self.x
        return [dx, dy]
\end{lstlisting}

\begin{warningbox}{backward 依赖 forward 的缓存}
注意一个关键细节：\texttt{backward} 方法需要用到前向传播时的输入值（如乘法门中的 $x$ 和 $y$）。因此 \texttt{forward} 方法必须将这些值\textbf{缓存}起来（存为类的属性）。这就是为什么在 PyTorch 的 \texttt{autograd.Function} 中，\texttt{forward} 需要用 \texttt{ctx.save\_for\_backward()} 保存张量。
\end{warningbox}

实际上，PyTorch 已经预定义了大量常用操作的 forward 和 backward 实现。用户只需像搭积木一样组合这些操作，PyTorch 会自动管理计算图和梯度传播。这正是``高中生也能做深度学习项目''的原因。

\subsection{本章小结}

现代深度学习框架将反向传播的``基础设施''（计算图管理、梯度传播）自动化，用户只需为每个操作定义前向计算和局部梯度。常见操作（线性层、激活函数、损失函数等）都已预实现，使用者通常不需要手动编写 backward 方法。但理解底层原理对于调试和设计新操作至关重要。

%% ========== Part 8: 梯度检验 ==========

